\documentclass[11pt,a4paper]{article}
\usepackage[margin=22mm]{geometry}\usepackage{fontspec}\usepackage{graphicx}\usepackage{booktabs}\usepackage{xcolor}\usepackage{hyperref}\usepackage{longtable}\usepackage{array}\usepackage{fancyhdr}\usepackage{titling}\setlength{\droptitle}{8mm}
\setmainfont{Apple SD Gothic Neo}\hypersetup{colorlinks=true,linkcolor=blue,urlcolor=blue}\pagestyle{fancy}\fancyhf{}\rhead{RAI Risk Taxonomy 2.0}\lhead{Human Review Recovery}\cfoot{\thepage}\setlength{\parskip}{6pt}\setlength{\parindent}{0pt}
\begin{document}
\title{RAI Risk Taxonomy 2.0\\Second Human-Review Recovery Technical Report}\author{Human-review recovery release}\date{29 August 2026}\maketitle
\par\section*{Executive summary}
This report documents the row-level reapplication of the second-round human-review instructions to 798 baseline L4 cards, followed by user-approved semantic deduplication and mechanism-level splitting. Neither EM nor Hybrid EM was rerun. After 13 non-representative cards were retired across ten consolidation clusters, ten compound cards were reviewed, two umbrella cards were retired, and seven distinct risk mechanisms were created. The final release contains 783 L4 cards: 613 General, 77 Agentic, and 93 Physical. No card remains assigned to Others.
\par\section{Recovery principles}
Valid edits and user-approved consolidations in the baseline commit were retained. Only explicit second-round instructions to delete, merge, split, remap, or rewrite were applied through source\_row\_id lineage. The L3 master prevailed over conflicting prompts. No L3 was created or modified, and its SHA-256 digest was frozen.
\par\section{Applied decisions}
\begin{table}[h]\centering\begin{tabular}{lr}\toprule Action & Reviewed rows\\\midrule Remap & 80\\Rewrite and retain & 39\\Split & 19\\Merge & 14\\Delete & 14\\\bottomrule\end{tabular}\caption{The 166 approved human-review decisions}\end{table}
\begin{figure}[h]\centering\includegraphics[width=.82\linewidth]{../figures/human_review_recovery_domain_counts.png}\caption{Final L4 cards by domain}\end{figure}
\begin{figure}[h]\centering\includegraphics[width=.88\linewidth]{../figures/human_review_recovery_actions.png}\caption{Distribution of approved human-review actions}\end{figure}
\par\section{Splitting compound cards}
Existing human-review split lineage was preserved while ten additional compound cards were reviewed. Where an existing card already represented a branch adequately, its unique meaning and source\_row\_id lineage were absorbed. Only seven independent harm mechanisms were created as new cards, and two overbroad umbrella cards were retired after their meanings were absorbed into specific cards. Every bilingual definition identifies an AI technical subject, a causal mechanism, and an adverse outcome.
\par\section{Semantic deduplication}
Cards were consolidated only when their risk event, affected interest, causal mechanism, and practical control objective were materially the same and any difference was limited to a channel, example, deployment context, or narrower restatement. Thirteen cards were retired across ten consolidation clusters covering self-harm and suicide facilitation, mass surveillance, contextual privacy disclosure, unauthorised use of copyrighted works, role-instruction safeguard bypass, evaluation deception, emotional dependence, socioeconomic inequality, accountability gaps, and domestic human-robot safety failure. Distinct protected characteristics and materially different harm mechanisms remained separate. The revised canonical cards preserve unique meaning and source\_row\_id lineage from every contributor.
\par\section{Integrity and validation}
The L3 master SHA-256 remained \texttt{e9439ced64fb49c1496f1955013b5f038ecc7d271b9d6c9704f1e1bf6b0094df}. L4 identifiers were reissued continuously within each final L3. Duplicate identifiers, exact duplicate cards, unknown L3 assignments, hierarchy mismatches, blank bilingual fields, and definitions without an AI technical subject all returned zero failures. Two independent executions produced identical SHA-256 hashes for the three L4 CSV files.
\par\section{Deliverables}
The release contains one L1 CSV, one combined L1/L2/L3 CSV, and separate General, Agentic, and Physical L4 CSVs. The validation directory preserves the 808-row instruction register, the 166 approved decisions, lineage edges from source\_row\_id to final L4, deletion tombstones, the semantic-consolidation log, the application log, and the validation record.
\par\section{Conclusion}
This recovery implements human-review intent rather than a new similarity-based classification. No Others assignments remain and the L3 master is unchanged. Any future remapping must be run only after a separate explicit user instruction.
\end{document}